% Stile «accademico»: verde bosco e oro, font con grazie (Charter), striscia laterale. Colori e font: modificare solo qui.
\usetikzlibrary{calc}
\usepackage{charter}
\usefonttheme{serif}
\definecolor{principale}{HTML}{1E4D3A}
\definecolor{accento}{HTML}{B8860B}
\definecolor{sfondochiaro}{HTML}{F5F2E8}
\usecolortheme[named=principale]{structure}
\setbeamersize{text margin left=1.1cm,text margin right=0.9cm}
\setbeamercolor{frametitle}{fg=principale}
\setbeamerfont{frametitle}{series=\bfseries,size=\Large}
\setbeamertemplate{frametitle}{\nointerlineskip\vskip3pt
  \begin{beamercolorbox}[wd=\paperwidth,leftskip=1.1cm,rightskip=0.9cm]{frametitle}\insertframetitle\end{beamercolorbox}}
\setbeamertemplate{background}{\begin{tikzpicture}[remember picture,overlay]
  \fill[principale](current page.north west)rectangle([xshift=0.35cm]current page.south west);
  \fill[accento]($(current page.north west)+(0.35cm,0)$)rectangle([xshift=0.45cm]current page.south west);
  \end{tikzpicture}}
\setbeamertemplate{itemize items}[triangle]
\setbeamercolor{itemize item}{fg=accento}
\setbeamercolor{block title}{bg=principale,fg=white}
\setbeamercolor{block body}{bg=sfondochiaro}
\setbeamercolor{block title alerted}{bg=accento,fg=white}
\setbeamercolor{block body alerted}{bg=sfondochiaro}
\setbeamertemplate{footline}{%
  \hfill\small\color{principale}\insertframenumber/\inserttotalframenumber\hspace{1em}\vspace{4pt}}
